\documentclass[a4paper,12pt]{article}

\usepackage[T2A]{fontenc}
\usepackage[utf8]{inputenc}
\usepackage[russian]{babel}
\usepackage{amsmath,amssymb,mathtools}
\usepackage{PTSans}
\renewcommand{\familydefault}{\sfdefault}
\usepackage{icomma}
\usepackage[a4paper,margin=20mm,headheight=25pt]{geometry}
\usepackage{microtype}
\usepackage{tikz}
\usetikzlibrary{arrows.meta,positioning,shapes.geometric}
\usepackage{tabularx,booktabs}
\usepackage{enumitem}
\usepackage[hidelinks]{hyperref}
\usepackage{parskip}
\usepackage{fancyhdr}
\usepackage{setspace}
\usepackage{xcolor}

\definecolor{accent}{HTML}{2F6F68}
\definecolor{darkaccent}{HTML}{234D49}
\definecolor{textgray}{HTML}{303638}
\definecolor{softgray}{HTML}{6E777A}
\definecolor{linegray}{HTML}{D4D9D8}

\geometry{top=19mm,bottom=20mm,left=20mm,right=20mm}
\onehalfspacing
\setlength{\parindent}{0pt}
\setlength{\parskip}{0.55em}
\setlist[itemize]{leftmargin=1.6em,itemsep=0.25em,topsep=0.25em}
\setlist[enumerate]{leftmargin=1.8em,itemsep=0.45em,topsep=0.3em}
\renewcommand{\arraystretch}{1.35}

\pagestyle{fancy}
\fancyhf{}
\fancyhead[L]{\small\color{softgray}Степени: формулы и вычисления}
\fancyhead[R]{\small\color{softgray}30.09.26}
\fancyfoot[C]{\small\color{softgray}\thepage}
\renewcommand{\headrulewidth}{0.3pt}
\renewcommand{\headrule}{\hbox to\headwidth{\color{linegray}\leaders\hrule height \headrulewidth\hfill}}

\newcommand{\topic}[1]{%
  \vspace{0.6em}
  {\large\bfseries\color{darkaccent}#1}\par
  \vspace{0.15em}\color{linegray}\hrule\color{textgray}\vspace{0.35em}
}
\newcommand{\keyidea}[1]{\textbf{\color{darkaccent}#1}}

\tikzset{
  flowstart/.style={ellipse,draw=darkaccent,thick,minimum width=36mm,minimum height=8mm,align=center,fill=white,font=\small},
  flowaction/.style={rectangle,rounded corners=2mm,draw=accent,thick,text width=106mm,minimum height=9mm,align=center,fill=white,font=\small},
  flowdecision/.style={diamond,aspect=3.25,draw=darkaccent,thick,text width=31mm,inner sep=1.6pt,align=center,fill=white,font=\small},
  flowarrow/.style={-{Stealth[length=3mm,width=2mm]},thick,draw=softgray}
}

\begin{document}
\color{textgray}

% -------------------- TITLE --------------------
\thispagestyle{empty}
\begin{center}
\vspace*{24mm}

{\Large\color{softgray}Чек-лист}\par
\vspace{7mm}
{\fontsize{27}{32}\selectfont\bfseries\color{darkaccent}Степени: формулы,\par приведение к общему основанию\par и вычисления}\par
\vspace{10mm}
{\large Марк Гукин}\par
\vspace{3mm}
{\normalsize 30 сентября 2026 года}\par

\vfill
{\small\color{softgray}Лёвин Артём Александрович эксклюзивно для Марка Гукина}\par
\vspace{10mm}

\begin{tikzpicture}[x=1cm,y=1cm]
  \draw[accent,line width=1.1pt]
    (-7.6,0) .. controls (-4.2,1.1) and (-1.6,-0.75) .. (1.3,0.25)
    .. controls (3.7,1.05) and (5.2,0.9) .. (7.6,0.1);
  \draw[linegray,line width=0.45pt]
    (-7.6,-0.38) .. controls (-3.6,0.45) and (0.8,-0.55) .. (7.6,-0.28);
\end{tikzpicture}
\end{center}
\clearpage

% -------------------- CORE CHECKLIST --------------------
\section*{Что нужно уметь после занятия}

\begin{itemize}[label={\(\square\)}]
  \item узнавать, когда выражение выгодно привести к одному основанию;
  \item раскладывать составные числа на простые множители и записывать их как степени;
  \item переводить корни в степени с дробным показателем;
  \item переводить дроби и десятичные дроби в степени с отрицательным показателем;
  \item применять свойства степеней в обоих направлениях;
  \item отличать \emph{основание степени} от \emph{показателя степени};
  \item после преобразований объединять одинаковые основания и сокращать выражение.
\end{itemize}

\topic{1. Формулы, которые должны узнаваться сразу}

Во всех примерах и упражнениях этого чек-листа основания положительны: \(a>0\), \(b>0\). Для корней считаем \(k\in\mathbb{N}\), \(k\ge2\). В этих условиях свойства можно применять без неоднозначности, связанной с дробными показателями:

\begin{align*}
 a^m\cdot a^n &= a^{m+n}, &
 \frac{a^m}{a^n} &= a^{m-n},\\[0.35em]
 (a^m)^n &= a^{mn}, &
 a^n b^n &= (ab)^n,\\[0.35em]
 \frac{a^n}{b^n} &= \left(\frac{a}{b}\right)^n, &
 \sqrt[k]{a^m} &= a^{m/k},\\[0.35em]
 a^{-n} &= \frac{1}{a^n}, &
 \left(\frac{a}{b}\right)^{-n} &= \left(\frac{b}{a}\right)^n.
\end{align*}

Дополнительно:
\[
 a^0=1,
 \qquad
 a^1=a.
\]

\keyidea{Самопроверка.} Формула применяется только тогда, когда её структура действительно совпадает со структурой выражения. Например, показатели складывают при умножении степеней с \emph{одинаковым основанием}.

\clearpage
\topic{2. Главная концепция: сначала привести основания к удобному виду}

На занятии использовалась фиксированная последовательность преобразований:

\begin{enumerate}
  \item \textbf{Составные числа вне показателей.} Если число является основанием степени, стоит под корнем или находится в знаменателе, разложите его на простые множители и по возможности запишите через степени простых чисел:
  \[
    32=2^5,\qquad 54=2\cdot3^3,\qquad 81=3^4,\qquad 25=5^2.
  \]
  \item \textbf{Корни вне показателей.} После разложения подкоренного числа переведите корень в дробную степень:
  \[
    \sqrt8=\sqrt{2^3}=2^{3/2},\qquad
    \sqrt[5]{81}=\sqrt[5]{3^4}=3^{4/5}.
  \]
  \item \textbf{Дроби.} Представьте знаменатель как степень и используйте отрицательный показатель:
  \[
    \frac{1}{2^3}=2^{-3},\qquad \frac18=\frac1{2^3}=2^{-3}.
  \]
  \item \textbf{Десятичные дроби.} Сначала переведите число в обыкновенную дробь, затем --- в степенную форму:
  \[
    0{,}2=\frac{2}{10}=\frac15=5^{-1},\qquad
    0{,}25=\frac14=2^{-2}.
  \]
\end{enumerate}

\keyidea{Концепция задаёт цель, формулы дают инструменты.} Сначала приводите выражение к небольшому числу удобных оснований, затем применяйте свойства степеней и сокращайте.

\topic{3. Что преобразуем, а что оставляем}

Если выражение имеет вид
\[
2^{\,3\sqrt7-1}\cdot 8^{\,1-\sqrt7},
\]
то \(8\) --- составное \emph{основание}, поэтому
\[
8^{1-\sqrt7}=(2^3)^{1-\sqrt7}=2^{3(1-\sqrt7)}.
\]
При этом выражение \(1-\sqrt7\) находится \emph{в показателе}. Сам корень \(\sqrt7\) здесь никуда переводить не требуется.

\keyidea{Критическая мысль.} Каскад «составные числа \(\to\) корни \(\to\) дроби» применяется к числам \emph{вне показателей}. Сложный показатель может оставаться сложным; после приведения оснований его достаточно корректно сложить, вычесть или умножить по формуле.

\clearpage

% -------------------- FLOWCHART --------------------
\thispagestyle{plain}
\begin{center}
{\Large\bfseries\color{darkaccent}Алгоритм: как упрощать выражение со степенями}
\end{center}
\vspace{5mm}

\begin{center}
\begin{tikzpicture}[node distance=6.2mm and 18mm]
  \node[flowstart] (start) {Получено выражение};
  \node[flowdecision,below=of start] (compound) {Есть составные числа вне показателей?};
  \node[flowaction,below=of compound] (factor) {Разложить их на простые множители;\\по возможности записать через степени простых чисел};
  \node[flowdecision,below=of factor] (roots) {Есть корни вне показателей?};
  \node[flowaction,below=of roots] (rootact) {Перевести корни в дробные степени:\\$\sqrt[k]{a^m}=a^{m/k}$};
  \node[flowdecision,below=of rootact] (fractions) {Есть дроби или десятичные числа вне показателей?};
  \node[flowaction,below=of fractions] (fractact) {Перевести их в удобную степенную форму;\\при необходимости применить отрицательный показатель};
  \node[flowaction,below=of fractact] (formulas) {Применить свойства степеней:\\раскрыть степень степени;\\объединить одинаковые основания;\\сложить или вычесть показатели};
  \node[flowstart,below=of formulas] (finish) {Вычислить результат и проверить преобразования};

  \draw[flowarrow] (start) -- (compound);
  \draw[flowarrow] (compound) -- node[right,font=\small]{да} (factor);
  \draw[flowarrow] (factor) -- (roots);
  \draw[flowarrow] (roots) -- node[right,font=\small]{да} (rootact);
  \draw[flowarrow] (rootact) -- (fractions);
  \draw[flowarrow] (fractions) -- node[right,font=\small]{да} (fractact);
  \draw[flowarrow] (fractact) -- (formulas);
  \draw[flowarrow] (formulas) -- (finish);

  % Ветви «нет» вынесены влево, чтобы стрелки не накладывались друг на друга.
  \draw[flowarrow] (compound.west) -- ++(-27mm,0) node[pos=0.28,above,font=\small]{нет} |- (roots.west);
  \draw[flowarrow] (roots.west) -- ++(-21mm,0) node[pos=0.28,above,font=\small]{нет} |- (fractions.west);
  \draw[flowarrow] (fractions.west) -- ++(-15mm,0) node[pos=0.28,above,font=\small]{нет} |- (formulas.west);
\end{tikzpicture}
\end{center}
\clearpage

% -------------------- EXAMPLES --------------------
\section*{Разобранные приёмы}

\topic{Пример 1. Разложение составного основания}
\[
\frac{2^3\cdot5^2}{10^2}
=\frac{2^3\cdot5^2}{(2\cdot5)^2}
=\frac{2^3\cdot5^2}{2^2\cdot5^2}
=2^{3-2}=2.
\]

Здесь всё подчинено одной цели: представить \(10\) через те же простые основания \(2\) и \(5\), после чего сократить степени.

\topic{Пример 2. Десятичные показатели и составное основание}
\[
5^{0{,}36}\cdot25^{0{,}32}
=5^{0{,}36}\cdot(5^2)^{0{,}32}
=5^{0{,}36+0{,}64}
=5.
\]

\topic{Пример 3. Корни как дробные показатели}
\[
\left(
\frac{2^{1/3}\cdot2^{1/4}}{\sqrt[12]{2}}
\right)^2
=\left(2^{\frac13+\frac14-\frac1{12}}\right)^2
=\left(2^{1/2}\right)^2
=2.
\]

Общий знаменатель показателей:
\[
\frac13+\frac14-\frac1{12}
=\frac4{12}+\frac3{12}-\frac1{12}
=\frac12.
\]

\topic{Пример 4. Корни внутри показателей сокращаются алгебраически}
\[
3^{\sqrt5+10}\cdot3^{-5-\sqrt5}
=3^{\sqrt5+10-5-\sqrt5}
=3^5=243.
\]

Здесь \(\sqrt5\) не требуется преобразовывать: основания уже совпадают, поэтому достаточно сложить показатели.

\topic{Пример 5. Составное основание при сложном показателе}
\begin{align*}
2^{3\sqrt7-1}\cdot8^{1-\sqrt7}
&=2^{3\sqrt7-1}\cdot(2^3)^{1-\sqrt7}\\
&=2^{3\sqrt7-1+3-3\sqrt7}\\
&=2^2=4.
\end{align*}

\clearpage
\section*{Типичные ошибки и финальная самопроверка}
\begin{itemize}[label={\(\square\)}]
  \item Складывать показатели при разных основаниях: \(2^a\cdot3^b\) нельзя объединить в одну степень, если показатели не совпадают и нет другой подходящей структуры.
  \item Забывать раскрывать «степень степени»: \((a^m)^n=a^{mn}\), показатели \emph{перемножаются}.
  \item Пытаться «избавиться от корня», который находится только в показателе степени, хотя основания уже совпали.
  \item Сразу считать большие числа вместо разложения основания на простые множители.
  \item Терять знак при делении степеней: \(a^m/a^n=a^{m-n}\).
  \item Использовать \(a^0=1\) при \(a=0\): выражение \(0^0\) в школьных преобразованиях не определяют.
\end{itemize}

\topic{Пять вопросов перед финальным ответом}
\begin{enumerate}
  \item Можно ли ещё разложить какое-либо составное основание?
  \item Остался ли корень вне показателя, который удобнее записать дробной степенью?
  \item Все ли одинаковые основания уже объединены?
  \item При раскрытии \((a^m)^n\) показатели точно перемножены?
  \item Итоговое выражение действительно проще исходного и все вычисления проверены?
\end{enumerate}

\keyidea{Ориентир из практики занятия.} Если после преобразований выражение не стало компактнее, полезно вернуться к основаниям и проверить, не пропущено ли разложение составного числа или возможность применить одну из формул степеней.

\clearpage

% -------------------- TRAINING --------------------
\section*{Тренировочный блок --- Путь А}
\textit{10 заданий. Решения не записаны: цель --- самостоятельно применить каскад «составные числа \(\to\) корни \(\to\) дроби \(\to\) свойства степеней».}

\begin{enumerate}
  \item Представьте в виде степени простого числа:
  \[
  64,\qquad 125,\qquad \frac1{27},\qquad 0{,}04.
  \]

  \item Представьте в степенном виде:
  \[
  \sqrt[4]{32},\qquad \sqrt[3]{25},\qquad \sqrt[6]{49}.
  \]

  \item Вычислите:
  \[
  \frac{2^5\cdot5^3}{10^3}.
  \]

  \item Вычислите:
  \[
  3^{0{,}45}\cdot9^{0{,}275}.
  \]

  \item Вычислите:
  \[
  \left(\frac{5^{1/2}\cdot5^{1/3}}{\sqrt[6]{5}}\right)^3.
  \]

  \item Вычислите:
  \[
  7^{\sqrt3+8}\cdot7^{-5-\sqrt3}.
  \]

  \item Вычислите:
  \[
  3^{2\sqrt2-2}\cdot27^{1-\frac23\sqrt2}.
  \]

  \item Вычислите:
  \[
  \frac{4^{3/2}\cdot8^{1/3}}{2^2}.
  \]

  \item Вычислите:
  \[
  \left(\frac{9^{1/4}\cdot27^{1/6}}{\sqrt3}\right)^4.
  \]

  \item Вычислите:
  \[
  16^{\sqrt5-1}\cdot8^{2-\frac43\sqrt5}\cdot2^{-1}.
  \]
\end{enumerate}

\clearpage

% -------------------- ANSWERS --------------------
\section*{Ответы}

\begin{tabularx}{\linewidth}{@{}>{\bfseries}p{12mm}X@{}}
\toprule
№ & Ответ \\
\midrule
1 & \(64=2^6;\quad 125=5^3;\quad \dfrac1{27}=3^{-3};\quad 0{,}04=\dfrac1{25}=5^{-2}\). \\
2 & \(2^{5/4};\quad 5^{2/3};\quad 7^{1/3}\). \\
3 & \(4\). \\
4 & \(3\). \\
5 & \(25\). \\
6 & \(343\). \\
7 & \(3\). \\
8 & \(4\). \\
9 & \(9\). \\
10 & \(2\). \\
\bottomrule
\end{tabularx}

\vfill
{\small\color{softgray}\textbf{Самопроверка после каждого задания:} все ли нужные составные числа вне показателей разложены; все ли корни вне показателей переведены в степени; одинаковы ли основания перед сложением показателей; правильно ли раскрыта степень степени; проверены ли арифметика и итоговый ответ.}

\end{document}
